\chapter{How to use this guide}\label{ch:howto}

This guide explains the order-slot manuscript of October 2026, one step at a time.
It assumes basic probability and no more. Chapter~\ref{ch:toolkit} supplies every other tool before it is needed.

Each chapter has the same parts:
\begin{itemize}
\item An \emph{In this chapter} box states the question and the plan.
\item \emph{In words} boxes give the intuition before and after each derivation.
\item Derivations show every step. A grey note at the right of each line names the rule that the step uses.
\item Boxes with a frame state a result of the paper with its exact hypotheses.
\item \emph{Watch out} boxes mark a hypothesis or a distinction that is easy to lose.
\item \emph{Check yourself} questions have full worked answers at the end of the chapter.
\item A \emph{What to remember} box closes the chapter.
\end{itemize}

There are no numerical examples in Chapters~\ref{ch:story} to~\ref{ch:empirical}.
The paper's example appears in symbols: two point-mass classes \(h\) and \(\ell\), a reserve \(p\), and an incumbent uniform on \([0,r]\).
Chapter~\ref{app:numbers} collects its values, for use with the figures.

Two passes work well.
\begin{enumerate}
\item First pass: read Chapter~\ref{ch:story}, then only the boxes of each later chapter.
\item Second pass: work through the derivations, and try each question before you read its answer.
\end{enumerate}

Chapters~\ref{ch:model} to~\ref{ch:ccc} carry the main argument. Chapters~\ref{ch:example} to~\ref{ch:empirical} build on it and can be read in any order.
